\section{Реализация шардированного ART-хранилища}

\subsection{Общая идея шардирования}

Одно из ограничений любой общей in-memory структуры заключается в
конкуренции за блокировки. Если все операции чтения и записи проходят
через одно дерево и один lock, то потоковая модель быстро упрётся в
contention. Поэтому в работе применяется шардирование.

Ключ префикса хешируется (используется FNV-32a), и значение хеша по
модулю числа шардов определяет, в какой шард он попадёт. Каждый шард
содержит собственное ART-дерево и собственный
\texttt{sync.RWMutex}. Такой подход даёт два преимущества:

\begin{itemize}
\item обновления разных префиксов, попадающих в разные шарды, не
конкурируют напрямую;
\item чтение и запись локализуются внутри одного шарда.
\end{itemize}

Устройство шардированного хранилища изображено на
рис.~\ref{fig:shards}.

\begin{figure}[ht]
\centering
\begin{tikzpicture}[node distance=0.4cm and 0.4cm,every node/.style={font=\footnotesize}]
\node[block, text width=3.0cm] (fv)  {Полный снимок\\ или update-событие};
\node[block, text width=4.0cm, below=of fv] (h)  {FNV-32a хеш\\ и маршрутизация по шардам};

\node[block2, below=1.0cm of h, xshift=-3.6cm] (lo0) {Загрузчик 0};
\node[block2, below=1.0cm of h, xshift=-1.2cm] (lo1) {Загрузчик 1};
\node[font=\large, below=1.0cm of h, xshift=0.5cm] {$\cdots$};
\node[block2, below=1.0cm of h, xshift=2.6cm] (lok) {Загрузчик $K{-}1$};

\node[block2, fill=gray!18, below=0.4cm of lo0] (a0) {ART$_0$};
\node[block2, fill=gray!18, below=0.4cm of lo1] (a1) {ART$_1$};
\node[font=\large, below=0.4cm of h, xshift=0.5cm, yshift=-1.0cm] {};
\node[block2, fill=gray!18, below=0.4cm of lok] (ak) {ART$_{K-1}$};

\draw[msarr] (fv) -- (h);
\draw[msarr] (h.south) -- (lo0.north);
\draw[msarr] (h.south) -- (lo1.north);
\draw[msarr] (h.south) -- (lok.north);
\draw[msarr] (lo0) -- (a0);
\draw[msarr] (lo1) -- (a1);
\draw[msarr] (lok) -- (ak);
\end{tikzpicture}
\caption{Устройство шардированного ART-хранилища}
\label{fig:shards}
\end{figure}

\subsection{Внутреннее представление ключа}

Для хранения в ART IP-префикс переводится в байтовое представление,
удобное для лексикографической навигации по дереву. Этот этап важен,
потому что ART работает на последовательности байтов, а не на
абстрактном объекте «префикс». Соответственно, корректное кодирование
ключа определяет корректность дальнейших операций поиска и обхода
\cite{leis2013art}. В реализации используется фиксированная длина
кодирования с явной длиной маски, что упрощает дальнейший longest
prefix match.

\subsection{Lookup по IP-адресу}

Точечный lookup для внешнего клиента реализует longest prefix match по
текущему состоянию хранилища. В зависимости от способа организации
ключей и навигации по дереву возможны разные стратегии:

\begin{itemize}
\item хранение всех префиксов и поиск наиболее специфичного кандидата;
\item последовательная проверка длин масок (до~33 для IPv4 и до~129
для IPv6);
\item более сложная навигация по radix-структуре с удержанием
последнего совпадения вдоль пути.
\end{itemize}

В текущей реализации lookup имеет приемлемые характеристики для IPv4
и более высокую стоимость для IPv6, что связано с особенностями
перебора масок и большей глубиной адресного пространства. Возможные
направления оптимизации обсуждаются в главе~10.

\subsection{Upsert и delete}

Главное достоинство нового backend-хранилища проявляется именно на
операциях обновления. \texttt{Upsert} затрагивает один шард и локально
модифицирует его дерево; \texttt{Delete} обрабатывает withdraw-событие
точно так же~-- локально, без перестройки остальных шардов. Это
принципиально отличается от поведения batch-структуры, в которой
каждое изменение маршрута приводило к полной перестройке таблицы.

\subsection{Dump полного состояния}

Полная выгрузка в ART-режиме реализуется как обход всех шардов и
сборка результатов. Это естественно делает \texttt{Dump} дороже, чем в
структуре, которая держит уже готовый линейный массив записей. По
сути, ART-подход платит за mutability дополнительной ценой обхода,
аллокаций и (опционально) последующей сортировки. В реализации
сортировка управляется параметром \texttt{ASMAP\_STORE\_ART\_DUMP\_SORT},
что позволяет в сценариях, где порядок не важен, экономить заметную
долю времени на больших таблицах.

Эта особенность не является дефектом реализации как таковым. Она
отражает фундаментальный компромисс: структура, оптимизированная под
частые инкрементальные обновления, редко оказывается одновременно
лучшей и для операции «отдать всё содержимое сразу».

\subsection{Контракт \texttt{prefixStore}}

С точки зрения исследовательской работы важна не сама конкретная
реализация, а наличие явного контракта \texttt{prefixStore}, общего
для двух backend-вариантов. Контракт задан минимальным набором
операций, описанных в табл.~\ref{tab:store-contract}; это делает
сравнение реализаций честным и сводит экспериментальный анализ к
сопоставлению поведения двух конкретных воплощений одного и того же
интерфейса.

\begin{table}[ht]
\centering
\caption{Контракт \texttt{prefixStore} и принципиальное поведение реализаций}
\label{tab:store-contract}
\small
\begin{tabular}{>{\raggedright\arraybackslash}p{2.6cm}>{\raggedright\arraybackslash}p{4.6cm}>{\raggedright\arraybackslash}p{5.6cm}}
\toprule
Операция & Реализация \textit{ranger} & Реализация \textit{шард. ART} \\
\midrule
\texttt{Replace} & полная замена слайса
                 & параллельная загрузка по шардам \\
\texttt{Upsert}  & пересборка всего слайса и полная замена ($O(n)$)
                 & локальная модификация одного шарда ($O(\log n / K)$ амортизированно) \\
\texttt{Delete}  & пересборка всего слайса
                 & локальное удаление в одном шарде \\
\texttt{Lookup}  & longest prefix match по линейной структуре
                 & longest prefix match по ART-дереву соответствующего шарда \\
\texttt{Dump}    & копирование готового слайса (1~аллокация)
                 & обход всех $K$ шардов, опционально сортировка \\
\bottomrule
\end{tabular}
\end{table}

\subsection{Целевая операция и интерпретация результатов}

Оценивать новое хранилище по одному только времени \texttt{Dump} или
по latency точечного \texttt{Lookup} некорректно. Целевая функция
работы другая -- сделать потоковую актуализацию реализуемой на
практике. Главной метрикой здесь становится стоимость
\texttt{Upsert} и \texttt{Delete}, а \texttt{Dump} и \texttt{Lookup}
играют роль «расходов на остальные сценарии». Эту установку важно
удерживать при чтении главы~9, где приведены эмпирические результаты.

\subsection{Согласованность параллельного доступа}

Шардирование само по себе не определяет модель согласованности; оно
лишь снижает контеншн между независимыми операциями. Чтобы
\texttt{Lookup} и \texttt{Upsert} безопасно сосуществовали внутри
одного шарда, в реализации используется стандартный
read-write-mutex. Все операции записи (\texttt{Upsert},
\texttt{Delete}, локальная фаза \texttt{Replace}) удерживают
write-lock соответствующего шарда; \texttt{Lookup} удерживает
read-lock того шарда, в котором происходит поиск.

Такая дисциплина даёт следующие гарантии:

\begin{itemize}
\item внутри одного шарда параллельные \texttt{Lookup}
выполняются без взаимного блокирования;
\item модификация одного шарда не блокирует операции в других
шардах;
\item \texttt{Dump} не требует глобального write-lock и реализуется
последовательным захватом read-lock каждого шарда либо
параллельным обходом независимых шардов с локальными read-lock.
\end{itemize}

Важно отметить, что согласованность \texttt{Dump} в этой модели
является \textit{шардовой}, а не глобальной snapshot-согласованностью
во всём дереве: каждое значение, появившееся в дампе, корректно
относительно своего шарда в момент чтения. Для прикладного сценария
получения полной таблицы это допустимо, поскольку клиент
интерпретирует результат как «текущее представление сервиса»,
а не как атомарный одномоментный срез всех шардов.

При необходимости более сильной семантики можно либо ввести
дополнительный механизм глобального snapshot (например, версионирование
шардов), либо ограничиться тем, что глобальный «строгий» snapshot
по построению не нужен в логике downstream-клиентов: они и так
строят своё состояние как $S_0 + \Delta$ через подписку.

\subsection{Поведение во время ресинхронизации}

Отдельный важный режим работы -- ресинхронизация. В этот момент
сервис должен:

\begin{enumerate}
\item получить новый снимок;
\item подготовить из него корректное внутреннее представление;
\item заменить текущее состояние store на новое;
\item возобновить обработку потока обновлений.
\end{enumerate}

В шардированной реализации эта последовательность выполняется как
скоординированная замена содержимого всех шардов; параллельный
\texttt{Lookup}, попадающий в момент перестроения, либо обслуживается
из частично уже обновлённого состояния, либо ожидает завершения
write-фазы конкретного шарда. Аномалии «исчезновения» ранее
существовавшего префикса в этот момент возможны, но они являются
ожидаемым следствием ресинхронизации и кратковременны. Учитывая, что
ресинхронизация происходит редко (только при сбоях или после
длительных пауз), такая модель является разумным инженерным
компромиссом.

%==============================================================
